% explainer-source-hash: e86e98979cf8c7389d590088e201dc0e7817163aa2a1f304cecfbf7c32a770ce
% explainer-source-commit: 3f1d18ace6744e68643813b24e0dc406bd4e6c7c
% explainer-generated: 2026-07-16
% Regenerate with /explain-all (see WIRING.md). Do not edit this stamp by hand:
% it is how the toolchain knows whether this document still matches the code.
\documentclass[11pt,a4paper]{article}
\input{preamble}

\title{\textbf{Voice Appointment Agent}\\[0.3em]
\large A Brief: the design, the evidence, and the gaps}
\author{Portfolio explainer, built from the repository source}
\date{}

\begin{document}
\maketitle
\thispagestyle{empty}

\section{The project in one page}

A small business loses money every time nobody answers the phone. This project is a working
answer: a real telephone number a person can speak to in ordinary sentences, which books,
reschedules, and cancels appointments in a real calendar.

Three off-the-shelf pieces, one piece of original work:

\begin{itemize}
  \item \textbf{Twilio} rents the phone line, answers the call, converts speech to text, and
  speaks replies back.
  \item \textbf{Gemini 2.5 Flash} reads each caller utterance and returns structured data:
  what they want, and the details extracted so far.
  \item \textbf{SQLite} stores appointments and refuses to double-book.
  \item \textbf{The original work} is the state machine in between, plus the architecture that
  makes an LLM-driven phone system testable without spending money on every run.
\end{itemize}

\begin{keyidea}{The claim to make, and its exact boundary}
The booking logic is thoroughly and genuinely verified: 47 free automated tests and 25 scripted
conversation scenarios, all passing, all reproduced independently while writing this document.

What is \emph{not} verified is that this survives a real human voice. The language model has
only ever been shown clean, typed English; Twilio's speech-to-text output is noisier. That gap
is real, and the project's README states it plainly rather than glossing it.
\end{keyidea}

\section{Architecture}

\begin{figure}[H]
\centering
\resizebox{0.96\textwidth}{!}{
\begin{tikzpicture}[node distance=10mm]
  \node[extbox, text width=20mm] (caller) {Caller\\{\scriptsize on a phone}};
  \node[extbox, right=26mm of caller, text width=24mm] (twilio) {\textbf{Twilio}\\{\scriptsize speech $\leftrightarrow$ text}};
  \node[box, right=26mm of twilio, text width=30mm] (app) {\textbf{app.py}\\{\scriptsize FastAPI: answers}\\{\scriptsize Twilio's webhooks}\\{\scriptsize with TwiML}};
  \node[box, right=26mm of app, text width=32mm] (conv) {\textbf{conversation.py}\\{\scriptsize the state machine:}\\{\scriptsize ALL business logic}};
  \node[box, below left=16mm and 2mm of conv] (ie) {\code{IntentEngine}\\{\scriptsize\itshape interface}};
  \node[box, below right=16mm and 2mm of conv] (cs) {\code{CalendarStore}\\{\scriptsize\itshape interface}};
  \node[extbox, below=12mm of ie] (gem) {\textbf{Gemini 2.5 Flash}\\{\scriptsize billed per call}};
  \node[store, below=12mm of cs] (db) {\textbf{SQLite}\\{\scriptsize appointments}};
  \node[store, above=13mm of conv] (sess) {\textbf{session\_store}\\{\scriptsize per-call memory}};

  % forward path: labels above the line
  \draw[flow] (caller) -- node[lbl,above]{voice} (twilio);
  \draw[flow] (twilio) -- node[lbl,above]{HTTP POST\\{\scriptsize per turn}} (app);
  \draw[flow] (app) -- node[lbl,above]{one\\utterance} (conv);
  \draw[flow] (conv) -- (ie);
  \draw[flow] (conv) -- (cs);
  \draw[flow] (ie) -- (gem);
  \draw[flow] (cs) -- (db);
  \draw[flow] (conv) -- node[lbl,right]{save} (sess);
  % return path: arcs below the row, clear of the boxes
  \draw[dflow] (conv) to[bend left=32] node[lbl,below]{reply text} (app);
  \draw[dflow] (app) to[bend left=32] node[lbl,below]{TwiML} (twilio);
  \draw[dflow] (twilio) to[bend left=32] node[lbl,below]{spoken reply} (caller);
\end{tikzpicture}}
\caption{\code{conversation.py} holds every business rule and depends on \emph{neither} Twilio
\emph{nor} Gemini. It talks only to the two interfaces. That is the whole design.}
\end{figure}

\subsection{The decision everything else rests on}

The conversation manager is handed its collaborators rather than creating them, and it only
ever knows them as \code{IntentEngine} and \code{CalendarStore}. So the identical state machine
runs in three worlds:

\begin{table}[H]
\centering
\small
\begin{tabular}{@{}llll@{}}
\toprule
\textbf{Where} & \textbf{Language model} & \textbf{Calendar} & \textbf{Cost per run} \\
\midrule
Production (uvicorn) & real Gemini, billed & SQLite file & real money \\
47 pytest tests & scripted double & temp SQLite & nothing \\
25 eval scenarios & scripted double & temp SQLite & nothing \\
6 marked tests & \textbf{real Gemini} & not involved & a fraction of a cent \\
\bottomrule
\end{tabular}
\caption{One state machine, four wirings. The expensive dependency is present exactly where it
must be proven, and absent everywhere else.}
\end{table}

\begin{keyidea}{Why this is the strongest thing in the project}
25 scenarios, several of them five turns long, would be over 100 billed model calls per run,
would take real time, and would be flaky, because language models are not deterministic.
Defining what the manager \emph{needs} rather than what it will \emph{use} converts an
unaffordable test suite into one that runs in under three seconds for free, and still exercises
the real state machine, the real database, and real conflict detection.

The discipline is what makes it credible: the fake engine lives in \file{eval/}, outside the
production package, and its own docstring says ``This is a TEST DOUBLE, not a production
fallback''. The shipped path always uses the real model, so the ``LLM-driven'' claim holds.
\end{keyidea}

\section{How a call works}

A phone call is not one conversation from the server's point of view; it is a series of
separate HTTP requests with no memory of each other. Everything the agent knows must therefore
be written to SQLite and read back every turn, keyed by Twilio's \code{CallSid}. That single
fact explains why \code{session\_store.py} exists at all.

\begin{figure}[H]
\centering
\resizebox{0.92\textwidth}{!}{
\begin{tikzpicture}[node distance=6mm]
  \node[box] (a) {Twilio POSTs \code{/voice/gather}\\{\scriptsize with the transcribed speech}};
  \node[box, below=5mm of a] (b) {load session + transcript from SQLite\\{\scriptsize the server had forgotten everything}};
  \node[box, below=5mm of b] (c) {\code{intent\_engine.parse\_turn(context)}\\{\scriptsize Gemini returns JSON: intent + slots + reply}};
  \node[box, below=5mm of c] (d) {state machine merges slots, applies the rules};
  \node[dec, below=6mm of d, text width=26mm] (e) {all slots filled\\AND confirmed?};
  \node[box, right=16mm of e, fill=softgrey] (f) {ask the next\\question};
  \node[box, below=7mm of e] (g) {validate: not past, within 9:00-17:00};
  \node[box, below=5mm of g] (h) {\code{calendar\_store.book(...)}\\{\scriptsize conflict check, then INSERT}};
  \node[box, below=5mm of h, fill=okgreen!12, draw=okgreen] (i) {reply with TwiML: \code{<Gather>} or \code{<Hangup>}};
  \draw[flow] (a)--(b); \draw[flow] (b)--(c); \draw[flow] (c)--(d); \draw[flow] (d)--(e);
  \draw[flow] (e) -- node[lbl,above]{no} (f);
  \draw[flow] (f.south) |- (i.east);
  \draw[flow] (e) -- node[lbl,left]{yes} (g);
  \draw[flow] (g)--(h); \draw[flow] (h)--(i);
\end{tikzpicture}}
\caption{One turn, end to end. Repeat until booked, hung up, or handed off.}
\end{figure}

\section{The engineering decisions worth defending}

\subsection{Turning the language model's flagship feature off}

Gemini 2.5 Flash ``thinks'' before answering by default. That helps hard reasoning and costs
seconds. On a phone call, several seconds of dead air is a product failure: the caller assumes
the line dropped. So the production configuration sets \code{thinking\_budget=0}, plus a strict
response schema and \code{temperature=0.1}.

\begin{keyidea}{The general lesson}
The right configuration depends on the product, not on the model's defaults. For a live phone
turn, latency beats depth. Knowing which clever feature to switch off is engineering; leaving
defaults on is not.

Equally: asking a model in prose to ``please reply in JSON'' sometimes yields ``Sure! Here's
the JSON:'' and then fails to parse. Supplying a formal \code{response\_schema} to the API
constrains the output structurally. That is the robust way to get machine-readable answers.
\end{keyidea}

\subsection{Reading the vendor documentation properly}

Twilio's \code{<Gather>} requests your server \emph{only when it captures speech}. On silence
it quietly moves on to the next instruction instead. The consequence would have been severe:
every ``sorry, could you repeat that?'' retry path would be unreachable for a silent caller,
who is precisely the caller who needs it. One attribute, \code{action\_on\_empty\_result=True},
routes silence back through the same state machine. This is the class of detail that separates
working telephony from a demo.

\subsection{Two safety nets, for two different failures}

\begin{itemize}
  \item \textbf{Retry cap (3 consecutive unclear turns)} catches the bad line and the noisy room.
  \item \textbf{Turn cap (14 total turns)} catches what the retry cap cannot: a caller whose
  every turn is perfectly intelligible but never completes anything. Say ``sometime in the
  morning'' forever and the retry counter resets every turn, because each turn was a valid
  continuation. Without the second cap the call never ends.
\end{itemize}

The test for the turn cap asserts \code{retry\_count == 0}, proving \emph{which} mechanism
saved the call rather than merely that something did. That is careful test design.

\subsection{Making time an input}

The conversation manager accepts an injectable clock (\code{now\_fn}) rather than calling
\code{datetime.now()}. Without this, a test asserting ``the appointment lands on 2026-07-15''
passes today and silently starts failing once that date passes, for reasons unrelated to any
code change. Every scenario pins time to \code{FIXED\_NOW = 2026-07-06 13:00}, chosen at
mid-afternoon so a same-day ``past'' time can be tested without also tripping the
business-hours rule.

\subsection{Half-open intervals}

Conflict detection is \code{a\_start < b\_end and b\_start < a\_end}. Strictly less-than, so an
appointment ending at 9:30 and one starting at 9:30 are back to back rather than in conflict,
which is what a busy clinic needs. Using \code{<=} would refuse a perfectly good booking. There
is a dedicated scenario for the boundary case.

\section{The evidence}

Every free number below was independently reproduced while writing this document, by installing
the pinned requirements into a clean virtual environment and running the suite. The 7 tests that
spend money were collected but \textbf{deliberately never executed}.

\begin{table}[H]
\centering
\small
\begin{tabular}{@{}lll@{}}
\toprule
\textbf{Claim} & \textbf{Observed} & \textbf{Verdict} \\
\midrule
54 tests collected & \code{54 tests collected} & confirmed \\
47 free, 7 billed & \code{47 passed, 7 deselected} & confirmed \\
25/25 scenarios pass & \code{25/25 scenarios passed} & confirmed \\
Clean venv install & exit code 0 & confirmed \\
Latency 1.8 s with thinking off & no artefact in repo & \emph{owner's record} \\
Live public-tunnel booking & no artefact in repo & \emph{owner's record} \\
\bottomrule
\end{tabular}
\end{table}

\subsection{What the 25 scenarios actually test}

\begin{keyidea}{Precisely, because the honest answer is narrower than it sounds}
\textbf{Genuinely exercised:} the real state machine, a real SQLite calendar on a real
temporary file, real conflict detection, real business-hours and past-date validation, real
retry and turn caps. Coverage spans booking (one-shot and multi-turn), rescheduling,
cancelling, conflicts, boundary-adjacent slots, ambiguous times and dates, unclear speech,
both safety caps, abandoned calls, and mid-call changes of mind.

\textbf{Substituted:} the language model, replaced by the regex-based scripted double.

\textbf{Not involved at all:} Twilio, FastAPI, HTTP, TwiML. Utterances are fed straight into
the manager. Those layers are covered separately by 7 endpoint tests.

So ``25/25'' is a strong statement about the booking \emph{logic} and no statement about
speech recognition. That framing is the one to offer first.
\end{keyidea}

\subsection{The tests were themselves tested}

A passing test proves nothing unless it can fail. The README claims the safety caps were
deliberately broken to confirm the scenarios noticed. That was re-verified here in a scratch
copy, leaving the repository untouched:

\begin{table}[H]
\centering
\small
\begin{tabular}{@{}lll@{}}
\toprule
\textbf{Sabotage} & \textbf{Result} & \textbf{Caught by} \\
\midrule
retry cap 3 $\rightarrow$ 999 & 24/25 & \code{unclear\_speech\_retry\_cap\_handoff} \\
turn cap 14 $\rightarrow$ 9999 & 24/25 & \code{max\_turn\_cap\_handoff} \\
\code{<} $\rightarrow$ \code{<=} in overlap & 24/25 & \code{boundary\_adjacent\_appointments\_no\_conflict} \\
\bottomrule
\end{tabular}
\caption{Each sabotage killed exactly the one scenario named for it, and nothing else. The
suite is load-bearing, not decorative.}
\end{table}

\begin{honest}{The scripted engine cannot fail to understand}
The fake engine was written to handle exactly the phrases the scenarios use; its docstring
concedes this. So no scenario can ever surface a misunderstanding. The 25 scenarios test the
state machine \emph{given perfect understanding}. Real understanding is tested only by the 6
marked Gemini tests, against clean typed English, and never against messy speech output.

This is the correct trade, and the project is open about it. Say it first: ``25/25 proves the
booking logic; I made the understanding layer a separate, smaller, honest experiment.''
\end{honest}

\section{Honest findings}

Found by reading the source and confirmed by running the real code in a scratch copy. None of
these appear in the README.

\begin{honest}{The webhook signature check fails open by default}
The most significant issue. Twilio signs each request so you can prove it came from Twilio.
The check is implemented correctly, but it is wrapped in \code{if base\_url and auth\_token:},
and \code{PUBLIC\_WEBHOOK\_BASE\_URL} is blank by default, with \file{.env.example} instructing
``leave blank for local development''.

Verified by driving the real app: with enforcement nominally \emph{on} and no base URL, a
request carrying \emph{no signature at all} returns \textbf{HTTP 200}. Set the base URL and the
same request correctly returns \textbf{403}.

There is a real reason (Twilio signs the public URL, so without it there is nothing to check
against). The problem is the direction of failure: a missing setting yields silently no
security rather than a refusal to start. Deploy behind a real URL and forget the variable, and
the booking endpoints are open to anyone. The fix is small: refuse to start, or make unsigned
requests an explicit opt-in.
\end{honest}

\begin{honest}{Phone numbers are stored in two incompatible formats, and a test locks it in}
Twilio gives caller ID as \code{+15551112222}; speech is parsed to bare digits
\code{5551112222}; lookup is an exact string match. Verified: an appointment booked after the
caller \emph{spoke} their number cannot be found when that person phones back. The caller who
was most helpful is the one who cannot cancel.

Sharpest detail: scenario 4 \emph{asserts this as correct}, checking that the caller-ID lookup
returns zero results. Its intent (a spoken number should override caller ID) is right; what is
missing is normalising both to one canonical format on input. A good example of a test encoding
an assumption rather than a requirement.
\end{honest}

\begin{honest}{The wrong appointment is targeted when a caller has two}
\code{list\_by\_phone} sorts furthest-future first, and the code takes the first match. Verified:
with appointments on 2026-08-01 and 2026-09-01 under one number, ``cancel my appointment''
targets \textbf{September}, keeping August. The agent does read the date back, so an attentive
caller can object, but ``my appointment'' in ordinary speech means the next one. No scenario
covers a caller with two active appointments, which is exactly why the suite stays green.
\end{honest}

\begin{honest}{The Google Calendar adapter has never run, and reading it shows why that matters}
Fully implemented against the real Calendar API, selected automatically if credentials are set,
but no credentials were ever provided. The README says so. Reading it reveals a concrete defect:
\code{cancel()} deletes the event and then fabricates a return value using
\code{datetime.now()}, and the conversation manager speaks that value back to the caller. On
this backend the agent would announce \emph{the current time} instead of the appointment's real
time. Invisible on SQLite, where \code{cancel()} returns the true row. That is the predictable
cost of shipping code you cannot run.
\end{honest}

\begin{honest}{Smaller items}
\begin{itemize}
  \item \code{find\_conflict} loads \emph{every} booked appointment ever made into Python on
  each booking, with no time filter and no index. Invisible at this scale; a one-line
  \code{WHERE} clause away from correct.
  \item \code{book()} checks for conflicts and inserts as two statements with no transaction, so
  two simultaneous bookings could both pass the check. One phone line makes it vanishingly
  unlikely, but ``why is your conflict check not atomic?'' is a fair question.
  \item \code{SessionState.stage} is declared, serialised into every row, and never read or
  written anywhere. Dead field from an earlier design.
  \item \code{cancel()} on an already-cancelled appointment silently succeeds, though the
  interface documents it as raising. Its sibling \code{reschedule()} does check. Inconsistent.
  \item No timezones anywhere; every time is one implicit local zone. Acknowledged in the README.
\end{itemize}
\end{honest}

\begin{honest}{Three claims rest on the owner's word}
The 1.8 second latency measurement, the 56 thinking-token figure, and the live public-tunnel
booking run all appear only in README prose; no log or artefact survives in the repository. The
tunnel run is the strongest verification claimed and was cleaned up afterwards by design.
Nothing suggests these are untrue, and the README's precision elsewhere argues the other way;
it even labels its own cost figure ``an estimate from per-token rates, not a verified line-item
invoice'', which is exactly right. But ``I ran it and here is the log'' and ``I ran it and
cleaned up'' are different, and an interviewer may probe. Next time, commit the log.
\end{honest}

\section{Security posture}

Secret hygiene is genuinely good. \file{.env} is a symbolic link pointing outside the
repository to a private secrets directory, so no credential is ever in the working tree.
The full git history (all 46 objects) was scanned for Twilio SIDs, Google API keys, private-key
blocks, and hard-coded tokens: clean. \file{.env.example} carries placeholders only. Every SQL
query is parameterised, which matters because a caller's spoken name reaches the database.
Configuration values are never logged.

The one real weakness is the fail-open signature check described above.

\section{Summary judgement}

\begin{keyidea}{Lead with these}
\begin{enumerate}
  \item \textbf{The interface seam.} Two small abstract classes turn an expensive,
  non-deterministic system into one with 72 free, deterministic checks. Best thing here.
  \item \textbf{Turning thinking mode off}, having measured it. Product engineering, not
  cargo-culted defaults.
  \item \textbf{The tests were tested.} Mutation testing reproduced exactly.
  \item \textbf{\code{actionOnEmptyResult}.} Evidence of reading the documentation properly.
  \item \textbf{The README's candour} about its own gaps exceeds most production documentation.
\end{enumerate}
\end{keyidea}

\begin{keyidea}{Raise these before someone else does}
\begin{enumerate}
  \item The signature check fails open without the base URL.
  \item Two phone formats, with a test that locks the bug in.
  \item The furthest-future appointment is targeted, not the soonest.
  \item ``25/25'' means the logic is right \emph{given perfect understanding}.
  \item No timezones; the Google adapter has never run; no real voice call has been made.
\end{enumerate}
\end{keyidea}

The honest one-line verdict: a well-architected demonstration of a real pattern, tested far
more rigorously than most portfolio projects, whose remaining gap (a real human voice on a real
line) is the one the author names first himself.

\vfill
\begin{center}
\small\itshape
Built from the repository source. Every claim is quoted from the code, verified by running it in
an isolated scratch copy, or explicitly labelled as the owner's own record. No phone call was
placed and no billed API test was executed in the making of this document.
\end{center}

\end{document}
